\subsubsection{sqrt-decomposition generico}
\begin{lstlisting}[style=mystyle, language=C++]
// TC: O(sqrt(N)) per operation
// usa: cuando quieres una EDD simple y lazy
#include <bits/stdc++.h>
using namespace std;

template<class T>
struct SqrtDecomp{
	int n, B;
	vector<T> a, bucket;
	SqrtDecomp(const vector<T>& v): a(v){
		n = v.size();
		B = (int)sqrt(n) + 1;
		bucket.assign(B, 0);
		for(int i=0;i<n;i++) bucket[i / B] += v[i];
	}
	void pointUpdate(int p, T val){
		bucket[p / B] += val - a[p];
		a[p] = val;
	}
	// range sum query
	T rangeSum(int l, int r){
		T res = 0;
		int bl = l / B, br = r / B;
		if(bl == br){
			for(int i=l;i<=r;i++) res += a[i];
		}else{
			for(int i=l;i<(bl+1)*B;i++) res += a[i];
			for(int i=bl+1;i<br;i++) res += bucket[i];
			for(int i=br*B;i<=r;i++) res += a[i];
		}
		return res;
	}
};

int main(){
	SqrtDecomp<int> sd({1, 2, 3, 4, 5, 6, 7, 8, 9, 10});
	cout << sd.rangeSum(2, 7) << "\n";  // 3+4+5+6+7+8 = 33
	sd.pointUpdate(4, 100);
	cout << sd.rangeSum(2, 7) << "\n";  // 128
	return 0;
}
\end{lstlisting}
